\newpage
\pagenumbering{arabic}
\setcounter{page}{1}
\pagestyle{mainstyle}
\thispagestyle{mainstyle}

\vspace*{0.5cm}
\begin{center}
\textbf{CHAPTER 1}\\[0.15cm]
\textbf{INTRODUCTION}
\end{center}
\label{chap:intro}
\vspace{0.4cm}

\subsection*{1.1 Background and Overview of Web3 Gaming}
Web3 applications introduce user-controlled non-custodial wallets, blockchain settlement, and cryptographic identity primitives into conventional web workflows. In modern digital wagering and prediction platforms, however, purely on-chain execution models introduce substantial block confirmation delays (such as 12-second block intervals on Ethereum Mainnet) and prohibitive gas transaction fees that conflict with fast-paced, interactive consumer gameplay. 

The CypherRoll platform resolves this fundamental tension through a state-of-the-art \textbf{3-Tier Hybrid Web3 Architecture}. In this design, high-frequency game execution and interactive 3D WebGL rendering are handled efficiently within the browser and application layer, while value custody and capital settlements are anchored to decentralized blockchain escrow contracts. This guarantees that players experience sub-second responsiveness without forfeiting decentralized custody or mathematical fairness.

\subsection*{1.2 Motivation of the Study}
The CypherRoll project is driven by four primary software engineering and architectural imperatives:
\begin{enumerate}[leftmargin=*]
    \item \textbf{Verifiable Mathematical Fairness:} Replacing centralized, opaque random number generators (``black-box RNGs'') with verifiable, deterministic commit-reveal cryptographic hashing protocols.
    \item \textbf{Low-Latency Interactive Gameplay:} Eliminating gas-fee frictions and transaction wait-times from individual gaming rounds by maintaining off-chain state execution backed by Three.js 60 FPS WebGL graphics.
    \item \textbf{Non-Password Cryptographic Wallet Identity:} Transitioning from vulnerable email/password accounts and intrusive KYC document harvesting to decentralized cryptographic wallet signatures via Sign-In with Ethereum (SIWE -- EIP-4361) and Solana Wallet Adapter.
    \item \textbf{Auditable Non-Custodial Capital Settlement:} Preventing operator solvency failures, commingling of player funds, and arbitrary withdrawal lockouts through smart contract escrow vaults deployed across EVM Layer-2 rollups.
\end{enumerate}

\subsection*{1.3 Problem Statement}
Traditional online gambling systems operate as centralized black boxes where server-side random number generation, player ledger databases, and treasury funds are held entirely by a single corporate operator. This centralized paradigm suffers from three critical industry-wide flaws:
\begin{itemize}[leftmargin=*]
    \item \textbf{Opaque and Unauditable Randomness:} Players have no mathematical mechanism to verify whether a dice roll, roulette spin, or multiplier crash was authentically random or algorithmically altered to maximize operator profit.
    \item \textbf{Custodial Counterparty Vulnerability:} Players must deposit funds directly into custodial operator accounts, leaving them completely vulnerable to insolvency, internal embezzlement, and frozen withdrawals.
    \item \textbf{Scalability Constraints of Pure On-Chain DApps:} Existing fully on-chain casinos execute every single wager as an individual blockchain transaction, resulting in exorbitant gas fees and multi-second latencies that make real-time gaming unviable.
\end{itemize}

\subsection*{1.4 Scope and Boundaries of the Project}
The scope of the CypherRoll project encompasses the end-to-end design, implementation, and empirical validation of a full-stack Web3 wagering application:
\begin{itemize}[leftmargin=*]
    \item Modern frontend application engineering using Next.js 14, TypeScript, Tailwind CSS, and Framer Motion.
    \item High-fidelity hardware-accelerated 3D WebGL rendering for \textbf{CypherDice} and \textbf{CypherCrash} built with Three.js and Draco-compressed GLB asset pipelines, complete with automatic canvas fallbacks for non-WebGL devices.
    \item Cryptographic Provably Fair Engine using HMAC-SHA256 with SHA-256 pre-committed server seeds, user-selected client seeds, and sequential nonces.
    \item Solidity Non-Custodial Smart Contract Escrow Vault (\texttt{CypherRollVault.sol}) deployed across EVM networks (Ethereum Sepolia and Base Sepolia testnets) with EIP-712 structured signature withdrawal authorization.
    \item Relational persistence and real-time state synchronization through Supabase PostgreSQL, incorporating atomic balance mutations, row-level security, and a live global betting feed.
    \item Automated AML/OFAC sanctions compliance screening, client-side PoW rate-limiting, and Kelly Criterion bankroll solvency controls.
\end{itemize}

\subsection*{1.5 Key Contributions of the System}
The primary contributions of this project include:
\begin{enumerate}[leftmargin=*]
    \item Designing and executing a provably fair cryptographic protocol that allows players to independently verify round outcomes in real-time or posthumously using open mathematical formulas.
    \item Developing a dual-chain Web3 authentication interface supporting EVM wallets (MetaMask, Coinbase Wallet, Rainbow) and Solana wallets (Phantom) under a unified session abstraction.
    \item Demonstrating an enterprise-grade Layer-2 settlement model achieving sub-cent gas fees and non-custodial escrow protection on Base Sepolia.
\end{enumerate}

\subsection*{1.6 Organization of the Report}
The remainder of this report is organized into ten chapters aligned with academic guidelines: Chapter 2 details the project objectives; Chapter 3 evaluates existing gaming architectures; Chapter 4 presents the proposed hybrid system; Chapters 5 through 8 provide the system requirements, architectural design, module breakdown, and implementation details; Chapter 9 presents testing and experimental results; and Chapter 10 concludes the report with future directions.
\label{chap:intro_end}
